\documentclass[11pt]{article}

% basic packages
\usepackage[margin=1in]{geometry}
\usepackage[pdftex]{graphicx}
\usepackage{amsmath,amssymb,amsthm}
\usepackage{custom}

% page formatting
\usepackage{fancyhdr}
\pagestyle{fancy}

\renewcommand{\sectionmark}[1]{\markright{\textsf{\arabic{section}. #1}}}
\renewcommand{\subsectionmark}[1]{}
\lhead{\textbf{\thepage} \ \ \nouppercase{\rightmark}}
\chead{}
\rhead{}
\lfoot{}
\cfoot{}
\rfoot{}
\setlength{\headheight}{14pt}

\linespread{1.03}
\setlength{\parindent}{0.2in}
\setcounter{secnumdepth}{2}
\setcounter{tocdepth}{2}

\begin{document}

% make title page
\thispagestyle{empty}
\bigskip \
\vspace{0.1cm}

\begin{center}
{\fontsize{22}{22} \selectfont Lecture Notes on}
\vskip 16pt
{\fontsize{36}{36} \selectfont \bf \sffamily
\href{https://fhsastrophysics.github.io/fhs-astrophysics/\#meeting-2627-1}{The Roman Space Telescope}}
\vskip 24pt
{\fontsize{18}{18} \selectfont \rmfamily Abir Mehta \& Dhruv Lagu}
\vskip 6pt
{\fontsize{14}{14} \selectfont \ttfamily amehta251@student.fuhsd.org}
\vskip 24pt
\end{center}

{\parindent0pt \baselineskip=15.5pt
These notes open the 2026--27 year. We first review the units astronomers use to
talk about distance, then turn to NASA's newest instrument, the Nancy Grace Roman
Space Telescope, which launched on August 30, 2026. We look at who Roman was, what
the telescope is built from, how it finds planets, and why it may finally pin down
dark energy.
}

\newpage
\microtoc
\newpage

% ==================================================
\section{How We Measure Space}

\subsection{Light-Year}

\begin{itemize}
\item A light-year is the distance light travels in one year, about
$5.88 \times 10^{12}$ miles.
\item It is a unit of \emph{distance}, not time.
\end{itemize}

\subsection{Astronomical Unit (AU)}

\begin{itemize}
\item 1 AU is the average distance between Earth and the Sun, about
$9.3 \times 10^{7}$ miles.
\item Useful for describing distances inside a planetary system.
\end{itemize}

\subsection{Parsec}

\begin{itemize}
\item 1 parsec $\approx 3.26$ light-years.
\item This is the unit professional astronomers actually use.
\end{itemize}

\subsection{For Scale}

\begin{itemize}
\item Alpha Centauri, the nearest star system, is about $4.35$ light-years away.
\item The Milky Way is roughly $100{,}000$ light-years across.
\end{itemize}

% ==================================================
\section{The Nancy Grace Roman Space Telescope}

\subsection{The Woman Behind the Telescope}

\begin{itemize}
\item Nancy Grace Roman (1925--2018) became NASA's first Chief of Astronomy in 1959
and its first woman executive.
\item She fought for years to get the Hubble Space Telescope funded and won, which is
why she is called the \textbf{Mother of Hubble}. Roman is her telescope.
\end{itemize}

\subsection{The Specs}

\begin{itemize}
\item \textbf{Launch.} Roman launched on August 30, 2026 at 4:26 AM PDT on a SpaceX
Falcon Heavy from Launch Complex 39A at Kennedy Space Center.
\item \textbf{Mirror.} A $2.4$ m primary mirror, the same size as Hubble's. NASA
received two such mirrors for free from a retired spy-satellite program.
\item \textbf{Camera.} A $300$-megapixel infrared camera built from $18$ detectors
in a curved arc. The shape became the mission logo.
\item \textbf{Location.} Parked at the second Lagrange point (L2), about
$1.5 \times 10^{6}$ km from Earth, right next to the James Webb Space Telescope.
\item \textbf{Data.} About $1.4$ TB per day, more in a single day than Hubble makes
in a whole year, over a planned $5$-year mission.
\end{itemize}

% ==================================================
\section{Hubble versus Roman}

\begin{itemize}
\item A single Roman image covers a patch of sky larger than the full Moon, at the
same sharpness as Hubble. A Hubble or JWST frame covers a patch about the size of a
grain of salt held at arm's length.
\item \textbf{Andromeda, side by side.} To image the Andromeda Galaxy, Hubble needed
$400$ individual images, $1{,}300$ imaging hours, and $4$ years. Roman needs $2$
images and about $1$ hour, and both results contain the same $1.5$ billion pixels.
\end{itemize}

% ==================================================
\section{What Roman Will Find}

\begin{itemize}
\item \textbf{Exoplanets.} About $100{,}000$ of them, roughly $16$ times more than
every exoplanet humanity has found so far, from this one mission.
\item \textbf{Rogue worlds.} Starless planets drifting alone through the galaxy, some
as small as Mars.
\item \textbf{Dark energy.} Dark energy makes up about $68\%$ of the universe and
drives its accelerating expansion, yet after $28$ years no one knows what it is.
Roman will map about $620$ million galaxies and $10{,}000$ exploding stars to trace
how space is stretching. Recent DESI results hint dark energy may be changing over
time, and Roman is expected to settle the question within $5$ years.
\end{itemize}

% ==================================================
\section{How Roman Finds Planets: Gravitational Microlensing}

\begin{itemize}
\item When a foreground star drifts in front of a more distant background star, the
foreground star's gravity bends and focuses the background star's light, briefly
brightening it. This is \textbf{gravitational microlensing}.
\item If the foreground star has a planet, the planet adds a small extra spike to the
light curve. That spike reveals worlds up to about $26{,}000$ light-years away.
\item Einstein predicted this effect in 1936 and thought it would never be observable
in practice. Roman turns it into a survey tool.
\end{itemize}

\end{document}
